\section{Transformation}
\label{transformation}
This section derives $H_d$ by transforming $H(t)$ and then expresses the master equation in Eq.~\eqref{Linblad} in the corresponding effective drive frame.

\subsection{Transformation of Hamiltonian}
In the regime $|g| \ll |\omega_c-\omega_b|\ll|\omega_c+\omega_b|$, counter-rotating terms can be neglected in the Hamiltonian $H_\text{sys}$ (Eq.\ \eqref{Linblad}), giving
\begin{equation}
\begin{split}
    H(t) \approx \,& \omega_b b^\dagger b + \omega_c c^\dagger c + \frac{K}{2} b^{\dagger 2} b^2+ g\,(bc^\dagger + b^\dagger c) + 2\Omega_0 \cos(\omega_d t)\,(b + b^\dagger).
\end{split}
\end{equation}
The static Hamiltonian $H_{\text{static}}$ consists of the first four terms. To remove the interaction term $g(bc^\dagger + b^\dagger c)$, a second-order Schrieffer--Wolff transformation is applied with $U_{\text{disp}} = e^{D_{\text{disp}}}$,
\begin{equation}
    H_{\text{static}}^{\text{disp}} = e^{-D_{\text{disp}}} H_{\text{static}} e^{D_{\text{disp}}}.
\end{equation}
The generator $D_{\text{disp}}$ is chosen to eliminate off-diagonal terms at first order,
\begin{equation}
\label{Ddisp}
    D_{\text{disp}} = \sum_n \sqrt{n+1}\,
    \frac{g}{\omega_b + nK - \omega_c}\,
    |n+1\rangle \langle n|\, c - \text{h.c.}
\end{equation}
Expanding in the detuning $\Delta_{\text{bc}}$ yields
\begin{equation}
    D_{\text{disp}} \approx
    \frac{g}{\Delta_{\text{bc}}} b^\dagger c
    - \frac{gK}{\Delta_{\text{bc}}^2} (b^\dagger b - 1)\, b^\dagger c
    + \mathcal{O}\!\left((nK/\Delta_{\text{bc}})^2\right).
\end{equation}

Applying the same transformation to the full (driven) Hamiltonian gives
\begin{align}
    H'_{\text{disp}}
    &= e^{-D_{\text{disp}}} H(t)\, e^{D_{\text{disp}}}
    \nonumber \\
    &\approx H_{\text{disp}}
    + \Omega_0 \cos(\omega_d t)\,
    \Bigl[e^{-D_{\text{disp}}} (b + b^\dagger) e^{D_{\text{disp}}}\Bigr]
    \nonumber \\
    &\approx H_{\text{disp}}
    + \Omega_0 \cos(\omega_d t)\,
    \Bigl[b^\dagger + [D_{\text{disp}}, b] + \text{h.c.}\Bigr],
\end{align}
where $H_{\text{disp}} \equiv e^{-D_{\text{disp}}} H_{\text{static}} e^{D_{\text{disp}}}$. Using the expansion above,
\begin{equation}
[D_{\text{disp}}, b] \approx
-\frac{g}{\Delta_{\text{bc}}}\, c
-\frac{Kg}{\Delta_{\text{bc}}^2}\, b\, c^\dagger
+\frac{2Kg}{\Delta_{\text{bc}}^2}\, b^\dagger b\, c .
\end{equation}
These terms rotate at frequencies set by detunings involving the $b$--$c$ manifold and the bare $c$ mode. Because $\omega_d$ is far detuned from these transitions, the rotating-wave approximation justifies neglecting them, leaving
\begin{equation}
H'_{\text{disp}} \approx H_{\text{disp}} + \Omega_0 \cos(\omega_d t)\,(b + b^\dagger).
\end{equation}

The explicit time dependence of the drive is removed by moving to the rotating frame defined by
$U_R = e^{i \omega_d t\, b^\dagger b}$. The transformed Hamiltonian is
\begin{align}
    H'_R
    &= U_R^\dagger \left( H'_{\text{disp}} - i \frac{\partial}{\partial t} \right) U_R
    \nonumber \\
    &\approx H_R + \left( \Omega_0 e^{-i \omega_d t} b + \text{h.c.} \right),
\end{align}
and a final RWA removes counter-rotating contributions at frequency $\pm2\omega_d$, yielding the effective Hamiltonian $H_R$.

The subsequent transformation with $U_d$ (Appendix~\ref{SWT}) maps $H_R$ to $H_d$.



\subsection{Transformation of $H_{\text{noise}}$}
\label{Hnoise_transfer}
We consider the flux-noise Hamiltonian
$H_{\text{noise}} = \frac{\partial \omega_b}{\partial \Phi} \delta \Phi(t) b^\dagger b$.
After applying the dispersive transformation $U_{\text{disp}} = e^{D_{\text{disp}}}$, the number operator becomes
\begin{align}
    e^{-D_{\text{disp}}} b^\dagger b e^{D_{\text{disp}}} &= e^{-D_{\text{disp}}} b^\dagger e^{D_{\text{disp}}} e^{-D_{\text{disp}}} b e^{D_{\text{disp}}} \nonumber \\
    &\approx \left( b^\dagger + [D_{\text{disp}}, b^\dagger] \right) \left( b + [D_{\text{disp}}, b] \right) \nonumber \\
    &\approx b^\dagger b + [D_{\text{disp}}, b^\dagger b] + [D_{\text{disp}}, b^\dagger][D_{\text{disp}}, b] \nonumber \\
    &\overset{\eqref{Ddisp}}{\approx} b^\dagger b + \frac{g^2}{\Delta_{bc}^2} c^\dagger c - \frac{4Kg^2}{\Delta_{bc}^3} b^\dagger b c^\dagger c \nonumber \\
    &\quad + \left( -\frac{g}{\Delta_{bc}} b c^\dagger + \frac{Kg}{\Delta_{bc}^2} b^\dagger b c^\dagger + \text{h.c.} \right).
\end{align}

The non-number-conserving contributions (e.g., terms proportional to $b c^\dagger$ and $b^\dagger b\,c^\dagger$, together with their Hermitian conjugates) drive transitions at frequencies near $\Delta_{bc}/2\pi\sim\mathrm{GHz}$. For $1/f$ flux noise, $S_{1/f}(\omega)$ is negligible at such frequencies, so these rapidly oscillating terms are discarded. The resulting dispersively transformed noise Hamiltonian is
\begin{equation}
    H_{\text{noise}}^{\text{disp}} \approx \frac{\partial \omega_b}{\partial \Phi} \delta \Phi \left( b^\dagger b + \frac{g^2}{\Delta_{bc}^2} c^\dagger c - \frac{4Kg^2}{\Delta_{bc}^3} b^\dagger b c^\dagger c \right).
\end{equation}

In the rotating frame defined by $U_R$, the remaining terms are diagonal in the relevant number operators, so the form is unchanged: $H_{\text{noise}}^R = H_{\text{noise}}^{\text{disp}}$.
Under the drive transformation $U_d = e^{D_d}$, the noise Hamiltonian is expanded to second order,
\begin{align}
    {H_{\text{noise}}^d}' &= U_d^\dagger H_{\text{noise}}^R U_d \nonumber \\
    &\approx H_{\text{noise}}^R + [D_d, H_{\text{noise}}^R] + \frac{1}{2} [D_d, [D_d, H_{\text{noise}}^R]] \nonumber ,
\end{align}
where we truncate the FTT Hilbert space to three levels in the regime of interest $|\Omega_0/K|\ll1$,
\begin{equation}
\begin{split}
    H_{\text{noise}}^R \approx \frac{\partial \omega_b}{\partial \Phi} \delta \Phi \bigg( & |e\rangle\langle e| + 2|f\rangle\langle f| + \frac{g^2}{\Delta^2_{bc}} c^\dagger c \\
    & - \frac{4Kg^2}{\Delta_{bc}^3} \left( |e\rangle\langle e| + 2|f\rangle\langle f| \right) c^\dagger c \bigg).
\end{split}
\end{equation}
Evaluating the required commutators:
\begin{align}
    &\left[ D_d, \, b^\dagger b\right] = -\sum_m \big[ \frac{\Omega_0}{\Delta_m} |g,m\rangle \langle e,m| \nonumber \\
    &+ \frac{\sqrt{2}\Omega_0}{K} |e,m\rangle \langle f,m| + \text{h.c.} \big] \nonumber\\
    &\frac{1}{2} \left[ D_d, \left[ D_d, b^\dagger b \right] \right]= \sum_m \Bigg[ 2 \left( \frac{\Omega_0}{\Delta_m} \right)^2 |g,m\rangle\langle g,m| \nonumber \\
    &- 2 \left( \frac{\Omega_0}{\Delta_m} \right)^2 |e,m\rangle\langle e,m|   - \frac{4\Omega_0^2}{K^2} |f,m\rangle\langle f,m| \Bigg]\nonumber \\
    &[D_d,c^\dagger c] = 0
\end{align}
gives
\begin{multline}
\label{noise Hamiltoniand}
    {H_{\text{noise}}^d}' \approx \frac{\partial \omega_b}{\partial \Phi} \delta \Phi \Bigg(  \frac{g^2}{\Delta_{bc}^2} c^\dagger c + \left(\frac{\Omega_0}{\Delta_m}\right)^2 |m\rangle\langle m| \\
    +\sigma^+ \sigma^-- \sigma^+ \sigma^- \left( \frac{4Kg^2}{\Delta_{bc}^3} c^\dagger c + 2\left(\frac{\Omega_0}{\Delta_m}\right)^2 |m\rangle\langle m| \right) \\
    - \sum_m \left[ \frac{\Omega_0}{\Delta_m} \sigma^- |m\rangle\langle m| + \text{h.c.} \right]\Bigg),
\end{multline}
with a two-level truncation in the FTT.

\subsection{Transformation of Lindbladian}

Following Ref.~\cite{Blais2021}, in the dispersive regime the reduced dynamics are described by the Lindblad master equation
\begin{align}
\label{Linblad}
    \dot{ \rho}_\textit{disp} &= -i[H_{\text{disp}},  \rho_\textit{disp}] + \gamma_{b\downarrow} \mathcal{D}[b] \rho_\textit{disp} + \gamma_{c\downarrow} \mathcal{D}[c]  \rho_\textit{disp}.
\end{align}
Here $\gamma_{b\downarrow}$ is the FTT relaxation rate, and $\gamma_{c\downarrow}=\frac{g^2}{\Delta_{bc}^2}\gamma_{b\downarrow}$ is the Purcell-induced cavity decay rate. The cavity relaxation is assumed to be Purcell-limited in state-of-the-art hardware designs \cite{longcoherence5}.

In the rotating frame defined by $U_R$, the equation retains the same Lindblad structure:
\begin{equation}
\label{LinbadR}
    \dot{\rho}_R = -i[H_R, \rho_R] + \gamma_{b\downarrow} \mathcal{D}[b]\rho_R + \gamma_{c\downarrow} \mathcal{D}[c]\rho_R,
\end{equation}
where ${\rho}_R = U_R^\dagger \rho_\textit{disp}U_R$.

Effective dissipation in the driven frame is obtained by transforming the jump operators with $U_d$. The FTT operator is first truncated to the three lowest levels ($|g\rangle, |e\rangle, |f\rangle$), so that $b \approx |g\rangle\langle e| + \sqrt{2}|e\rangle\langle f|$, and the drive transformation $U_d=e^{D_d}$ is applied. Projecting the result to the two-level FTT subspace then yields
\begin{align}
\label{transformb}
    e^{-D_d} b e^{D_d} &\approx b + [D_d,b] +\frac{1}{2}[D_d,[D_d,b]] \nonumber \\
    &\approx \sigma^- + \sum_m \Omega_0 \bigg( -\frac{1}{\Delta_m} |m\rangle\langle m| \nonumber \\
    &\quad + \left[ \frac{2}{\Delta_m} - \frac{2}{K} \right] \sigma^\dagger \sigma |m\rangle\langle m| \bigg) \nonumber \\
    &\quad + \sum_m \Omega_0^2 \bigg[ \left( \frac{1}{K \Delta_m} - \frac{1}{\Delta_m^2} \right) \sigma^+ |m\rangle\langle m| \nonumber \\
    &\quad + \left( \frac{2}{K \Delta_m} - \frac{1}{\Delta_m^2} \right) \sigma^- |m\rangle\langle m| \bigg],
\end{align}
where we truncate FTT to two levels.
Similarly, the cavity operator $c$ transforms under the influence of the drive:
\begin{align}
\label{transformc}
    &e^{-D_d} c e^{D_d}\nonumber \\
    &\approx c + \sum_m \frac{\sqrt{m+1} \Omega_0 \chi}{\Delta_{ m} \Delta_{ m+1}} \big( (\sigma^+ - \sigma^-)|m\rangle\langle m+1|    \nonumber \\
    &\quad + \sum_m \sqrt{m+1} \Omega_0^2 \left( \frac{1}{\Delta_{ m+1}^2} - \frac{1}{\Delta_{ m}^2} \right) |m\rangle\langle m+1| \nonumber \\
    &\quad - \sum_m \sqrt{m+1} \frac{\Omega_0^2 (\chi^2 + \Delta_{m} \chi)}{(\Delta_{m} \Delta_{ m+1})^2} \sigma^+ \sigma^- |m\rangle\langle m+1|
\end{align}

In this driven frame, the Lindblad master equation becomes
\begin{align}
\label{Linbdadappendix}
    \dot{\rho}_d &= -i[H_d, \rho_d] + \gamma_{b\downarrow} \mathcal{D}\left[ e^{-D_d} b e^{D_d} \right] \rho_d \nonumber \\
    &+ \gamma_{c\downarrow} \mathcal{D}\left[ e^{-D_d} c e^{D_d} \right] \rho_d \nonumber \\
    &\approx-i[H_d, \rho_d] +\nonumber \\
    &\gamma_{b\downarrow} \mathcal{D}\left[ \sigma^- + \sum_m \Omega_0^2 \left(\frac{2}{K\Delta_m} - \frac{1}{\Delta_m^2}\right) \sigma^- |m\rangle\langle m| \right] \rho_d \nonumber \\
    &\quad + \gamma_{b\downarrow} \mathcal{D}\left[ -\sum_m \frac{\Omega_0}{\Delta_m} |m\rangle\langle m| \right] \rho_d \nonumber \\
    &\quad + \gamma_{b\downarrow} \mathcal{D}\left[ \sum_m \Omega_0^2 \left(\frac{1}{K\Delta_m} - \frac{1}{\Delta_m^2}\right) \sigma^+ |m\rangle\langle m| \right] \rho_d \nonumber \\
    &\quad + \sum_m \gamma_{c\downarrow, m} \mathcal{D}\big[ |m\rangle\langle m+1| \big] \rho_d \nonumber  \\
    &+\mathcal{O}\left( \frac{g^2}{\Delta_{bc}^2} \frac{(m+1)\Omega_0^2 \chi^2}{(\Delta_m \Delta_{m+1})^2} \right),
\end{align}
where $\rho_d = U_d^\dagger \rho_R U_d$. The effective cavity decay rate $\gamma_{c\downarrow, m}$ is perturbatively modified by the drive:
\begin{equation}
    \gamma_{c\downarrow, m} \approx (m+1) \left( 1 - \Omega_0^2 \left[ \frac{1}{\Delta_{m+1}^2} - \frac{1}{\Delta_m^2} \right] \right) \gamma_{c\downarrow}.
\end{equation}

Equation \eqref{Linbdadappendix} uses the secular approximation, valid when the relevant detunings and mode frequencies exceed the decay rates ($\widetilde\Delta_\textit{bd}, \omega_c \gg \gamma_{\textit{b}\downarrow}$). In the interaction picture, the transformed jump operator decomposes into spectral components $\sum_k L_k$ oscillating at distinct frequencies $\omega_k$. When $\gamma \ll |\omega_j-\omega_k|$, rapidly oscillating cross-terms can be neglected, so that dissipation channels decouple:
\begin{equation}
\gamma \mathcal{D}\left[ \sum_k L_k \right] \rho \approx \sum_k \gamma \mathcal{D}[L_k] \rho.
\end{equation}

\subsection{Further approximation under two regimes}
\label{Linbladian regime}
We have transformed the Lindbladian in Eq.~\eqref{Linblad} into an effective drive frame, yielding the dissipators in Eq.~\eqref{Linbdadappendix} and the noise Hamiltonian in Eq.~\eqref{noise Hamiltoniand}. We then derive simplified expressions for the Lindbladian and $H_{\mathrm{noise}}^{d}$ in the selective and unselective regimes discussed in Sec.~\ref{subsec:ssconditiondiscussion}.


\paragraph{Unselective regime}
Retain up to second order of $\mathcal{O}((\Delta_{bd}/\chi)^2)$, dissipators in Eq.~\eqref{Linbdadappendix} and the noise Hamiltonian in Eq.~\eqref{noise Hamiltoniand} reduce to 
\begin{align}
\dot{\rho} &= -i [H_d + H_{\text{noise}}^d, \rho]  + \widetilde{\gamma}_{b\downarrow} \mathcal{D}\left[ \sigma^- (1 - \lambda_\downarrow c^\dagger c) \right] \rho \nonumber \\
&\quad + 2\gamma_{c\phi} \mathcal{D}\left[ c^\dagger c \right] \rho + \gamma_{b\uparrow} \mathcal{D}\left[ \sigma^+ (1 - \lambda_\uparrow c^\dagger c) \right] \rho + \widetilde{\gamma}_{c\downarrow} \mathcal{D}[c] \rho,
\end{align}
with the noise operator 
\begin{align}
H_{\text{noise}}^d &\approx \delta \omega_b \Bigg[ \left( \frac{g^2}{\Delta_{bc}^2} - \frac{2 \Omega_0^2 \chi}{\Delta_{bd}^3} \right) c^\dagger c \nonumber \\
&\quad - 2 \left( \frac{\Omega_0}{\Delta_{bd}} \right)^2 \sigma^+ \sigma^- + \left( \frac{4 \Omega_0^2 \chi}{\Delta_{bd}^3} - \frac{4 K g^2}{\Delta_{bc}^3} \right) \sigma^+ \sigma^- c^\dagger c \nonumber \\
&\quad - \frac{\Omega_0}{\Delta_{bd}} \sigma_x + \frac{\Omega_0 \chi}{\Delta_{bd}^2} \sigma_x c^\dagger c \Bigg],
\end{align}
and normalized FTT relaxation rate, pure cavity dephasing rate, FTT excitation rate, normalized cavity decay rate respectively:
\begin{align}
\widetilde{\gamma}_{b\downarrow} &\approx \gamma_{b\downarrow} \left[ 1 + 2\left( \frac{2\Omega_0^2}{K\Delta_{bd}} - \frac{\Omega_0^2}{\Delta_{bd}^2} \right) \right] +\mathcal{O}\bigg(\frac{\Omega^4_0}{\Delta^4_{bd}}\bigg) \\
\gamma_{c\phi} &= \gamma_{b\downarrow} \left( \frac{\Omega_0 \chi}{\Delta_{bd}^2} \right)^2/2 \\
\gamma_{b\uparrow} &= \gamma_{b\downarrow} \left[ \Omega_0^2 \left( \frac{1}{K\Delta_{bd}} - \frac{1}{\Delta_{bd}^2} \right) \right]^2  \\
\widetilde{\gamma}_{c\downarrow} &\approx \gamma_{c\downarrow} \left( 1 + \frac{2\chi \Omega_0^2}{\Delta_{bd}^3} +\mathcal{O}\bigg(\frac{\chi^2 }{\Delta^2_{bd}}\bigg)\right).
\end{align}
When small parameters
\begin{align}
\lambda_\downarrow &\approx \frac{2\chi}{\Delta_{bd}} \Omega_0^2 \left( \frac{1}{K\Delta_{bd}} - \frac{1}{\Delta_{bd}^2} \right) + \mathcal{O}\bigg(\frac{\Omega^4_0}{\Delta^4_{bd}}\bigg) \\
\lambda_\uparrow &= \frac{\chi}{\Delta_{bd}} \left( \frac{2K - \Delta_{bd}}{K - \Delta_{bd}} \right).
\end{align}
The cavity-photon dependence of the decoherence rates originates from the drive-induced nonlinearity. For higher cavity occupations, the AC-Stark-shifted detuning $\Delta_m$ increases, which suppresses the FTT heating rate and drives the relaxation rate toward $\widetilde{\gamma}_{b\downarrow}$. In the main text, we neglect this photon-number dependence to keep the discussion focused on the leading scaling. Using Eq.~\eqref{eq:ss_unselective}, we then recover the simplified master equation quoted in subsection~\ref{subsec:unselective-dephasing}.


\paragraph{Selective regime.}
In the selective regime, the condition $|\Delta_{m^\star}/\chi|\ll 1$ holds. Combining the relation $\chi \approx 2\frac{g^2}{\Delta_{bc}^2}K$ with the hierarchy of energy scales gives
\begin{equation}
|\Omega_0| \ll |\Delta_{m^\star}| \ll |\chi| \ll |K|.
\end{equation}
Accordingly, drive-induced corrections to the dissipators are retained up to $\mathcal{O}(\Omega_0/K)$. Terms of the form
$\Omega_0/(\Delta_{m^\star}+m\chi)$ in Eq.~\eqref{Linbdadappendix} scale as $\mathcal{O}(\Omega_0/\chi)$ for $m\neq m^*$ (since $|\Delta_{m^\star}+m\chi|\sim |m\chi|$), whereas the resonant contribution scales as $\Omega_0/\Delta_{m^\star}$ for $m=m^*$. To leading order of $\Omega_0/\Delta_{m^\star}$, the master equation reduces to
\begin{align}
&\dot{\rho} \approx -i\,[H_d + H_{\mathrm{noise}}^{d},\,\rho] \nonumber\\
&+ \gamma_{b\downarrow}\,\mathcal{D}\!\left[\sigma^- - \frac{\Omega_0^2}{\Delta_{m^\star}^2}\,\sigma^- |m^\star\rangle\langle m^\star|\right]\rho
+ 2\gamma_{c\phi}^{\mathrm{sel}}\,\mathcal{D}\!\left[|m^\star\rangle\langle m^\star|\right]\rho \nonumber\\
&+ \gamma_{b\uparrow}^{\mathrm{sel}}\,\mathcal{D}\!\left[\sigma^+ |m^\star\rangle\langle m^\star|\right]\rho
+ \sum_m \widetilde{\gamma}_{c\downarrow,m}^{\mathrm{sel}}\,\mathcal{D}\!\left[|m\rangle\langle m+1|\right]\rho ,
\end{align}
and the noise Hamiltonian is
\begin{align}
H_{\mathrm{noise}}^{d}
\approx \delta\omega_b(t)\Bigg[
&\frac{g^2}{\Delta_{bc}^2}\,c^\dagger c
+ \frac{g^2}{\Delta_{bc}^2}\,|m^\star\rangle\langle m^\star| \nonumber\\
&+ \sigma^+\sigma^-\!\left(\frac{4Kg^2}{\Delta_{bc}^3}\,c^\dagger c
+ 2\frac{g^2}{\Delta_{bc}^2}\,|m^\star\rangle\langle m^\star|\right) \nonumber\\
&- \frac{g}{\Delta_{bc}}\,\sigma^+ |m^\star\rangle\langle m^\star|
+ \mathrm{h.c.}
\Bigg].
\end{align}
Here, the selective-regime sweet-spot condition Eq.~\eqref{eq:Dm_selective} has been used. At this operating point, the longitudinal component of $H_{\mathrm{noise}}^{d}$ produces no dephasing between \(|g,m^\star\rangle\) and \(|g,m^\star+1\rangle\). Meanwhile residual cavity dephasing from AC-Stark-shift fluctuations induced by \(\delta\omega_b(t)\) is negligible for the dynamics of interest because the FTT remains near its ground state when initialized in \(|g\rangle\). The transverse component instead induces depolarization; within a Markovian approximation, the corresponding excitation rate due to \(1/f\) noise is
\(
\gamma_{b\uparrow}^{1/f,\mathrm{sel}}
= \left( \partial \omega_b / \partial \Phi \right)^2
\left( g / \Delta_{bc} \right)^2
S_{1/f}(\Delta_{bd}) .
\)

The dissipative processes in the effective drive frame can be summarized as follows:
(i) FTT relaxation with a cavity-state-dependent correction, together with cavity decay at a renormalized rate
\(
\gamma_{c\downarrow,m}^{\mathrm{sel}}
= \gamma_{c\downarrow}^{\mathrm{sel}}\left(1 + \delta_{m^*,m} g^2/\Delta_{bc}^2 \right),
\)
where
\(
\gamma_{c\downarrow}^{\mathrm{sel}} = (g^2/\Delta_{bc}^2)\gamma_{b\downarrow} 
\)
is the Purcell decay rate;
(ii) FTT heating at rate
\(
\gamma_{b\uparrow}^{\mathrm{sel}} = (g^4/\Delta_{bc}^4)\gamma_{b\downarrow}
\)
when the cavity occupies \(|m^\star\rangle\);
and (iii) pure dephasing between \(|m^\star\rangle\) and \(|m^\star+1\rangle\) with rate
\(
\gamma_{c\phi}^{\mathrm{sel}} = (g^2/\Delta_{bc}^2)\gamma_{b\downarrow}^{\mathrm{sel}}/2.
\)

Dephasing between \(|g,m^\star\rangle\) and \(|g,m^\star+1\rangle\) is therefore limited by cavity pure dephasing at rate \(\gamma_{c\phi}^{\mathrm{sel}}\propto g^2/\Delta_{bc}^2\), and by FTT heating at rate \(\gamma_{b\uparrow}^{\mathrm{sel}}\propto 1/\Delta_{bd}\) due to photon shot noise in the selective regime.

